\documentclass[british]{article}

\usepackage[margin=1.5in]{geometry}

\title{Local analytic geometry, II. Theory of finite morphisms}
\author{Christian Houzel}
\date{1961}

% Standard
\usepackage{amssymb,amsmath,amscd}
\usepackage{enumerate}
\usepackage{tikz-cd}
\usepackage{mathtools}
\usepackage{hyperref}
\usepackage{booktabs}


%Bibliography
\usepackage[backend=biber,maxnames=99,maxalphanames=4]{biblatex}
\addbibresource{SHC-13_2-19.bib}


% Colours
\usepackage{xcolor}
\hypersetup{colorlinks,linkcolor={purple!50!black},citecolor={purple!50!black},urlcolor={purple!80!black}}


% Fonts
\usepackage{mathrsfs}
\usepackage{fouriernc}


% Things for Quarto/pandoc
\definecolor{quarto-callout-color-frame}{HTML}{000000}
\usepackage{calc,array}
\usepackage{longtable}
\providecommand{\tightlist}{%
\setlength{\itemsep}{0pt}\setlength{\parskip}{0pt}}
\newcounter{none}


% Header and footer
\usepackage{fancyhdr}
\usepackage{lastpage}
\usepackage{datetime2}
\pagestyle{fancy}
\fancypagestyle{plain}{}
\fancyhf{}
\renewcommand{\headrulewidth}{0pt}
\cfoot{\small\thepage\ of \pageref*{LastPage}}
\lfoot{\footnotesize Build date: \today}


% Theorem environments
\usepackage{amsthm}
\newenvironment{itenv}[1]
  {\phantomsection\par\smallskip\noindent\textbf{#1.}\itshape}
  {\par\smallskip}
\newenvironment{rmenv}[1]
  {\phantomsection\par\smallskip\noindent\textbf{#1.}\rmfamily}
  {\par\smallskip}


% Shortcuts
\newcommand{\oldpage}[1]{\marginpar{\footnotesize$\Big\vert$ \textit{p.~#1}}}


%% Document %%

\usepackage{embedall}
\begin{document}

\maketitle
\thispagestyle{fancy}

\setcounter{footnote}{0}


%% Content %%

\emph{This document is a translation into English of the following:}

\begin{quote}
Houzel, C. ``Géométrie analytique locale, II. Théorie des morphismes
finis''. \emph{Séminaire Henri Cartan} \textbf{13 (2)} (1960--61), Talk
no. 19.
\href{https://www.numdam.org/item/SHC_1960-1961__13_2_A6_0}{numdam.org/item/SHC\_1960-1961\_\_13\_2\_A6\_0}
\end{quote}

\emph{The translator (\href{https://thosgood.net}{Tim Hosgood}) takes
full responsibility for any errors introduced in this document, and
claims no rights to any of the mathematical content.}

\tableofcontents

\providecommand{\scr}[1]{{\mathscr{#1}}}
\renewcommand{\cal}[1]{{\mathcal{#1}}}
\renewcommand{\frak}[1]{{\mathfrak{#1}}}
\renewcommand{\geq}{\geqslant}
\renewcommand{\leq}{\leqslant}

\providecommand{\Cat}[1]{\mathsf{#1}}

\providecommand{\sh}[1]{\scr{#1}}
\providecommand{\OO}{\sh{O}}
\providecommand{\op}{\mathrm{op}}

\providecommand{\AA}{\sh{A}}
\providecommand{\BB}{\sh{B}}
\providecommand{\MM}{\sh{M}}
\providecommand{\NN}{\sh{N}}

\providecommand{\Hom}{\operatorname{Hom}}
\providecommand{\shHom}{\underline{\Hom}}

\providecommand{\Specan}[1]{\operatorname{Spec}_\mathrm{an}(#1)}
\providecommand{\Alg}[1]{{#1}\text{-}\Cat{Alg}}
\providecommand{\Mod}[1]{{#1}\text{-}\Cat{Mod}}

\section{Analytic spectrum of an algebra of finite
presentation}\label{section-1}

\oldpage{19-01}

\protect\phantomsection\label{definition-1}
\begin{rmenv}{Definition 1}
Let \(S\) be a ringed space. We say that an \({\mathscr{O}}_S\)-algebra
\({\mathscr{A}}\) is of \emph{finite presentation} if every \(s\in S\)
admits a neighbourhood \(U\) such that \({\mathscr{A}}|U\) is isomorphic
to a sheaf of the form \[
  \frac{({\mathscr{O}}_S|U)[t_1,\ldots,t_n]}{(f_1,\ldots,f_m)}
\] with \(f_1,\ldots,f_m\in\Gamma(U,{\mathscr{O}}_S[t_1,\ldots,t_n])\).

\end{rmenv}

That is, \({\mathscr{A}}\) is locally generated over \({\mathscr{O}}_S\)
by a finite number of sections subject to a finite number of relations.

We are interested in the case where \(S\) is an analytic space, over a
complete valued non-discrete base field \(k\). For every analytic space
\(T\) over \(S\), consider the sheaf
\({\mathscr{A}}_T={\mathscr{A}}\times_S T\), the inverse image of
\({\mathscr{A}}\) on \(T\), and the set
\(F_{\mathscr{A}}(T)=\operatorname{Hom}_{{\mathscr{O}}_T}({\mathscr{A}}_T,{\mathscr{O}}_T)\)
of algebra homomorphisms; we immediately see that
\(T\mapsto F_{\mathscr{A}}(T)\) defines a contravariant functor
\(F_{\mathscr{A}}\colon(\mathsf{An})^\mathrm{op}/S\to\mathsf{Set}\).

\protect\phantomsection\label{proposition-1}
\begin{itenv}{Proposition }
Let \({\mathscr{A}}\) be an algebra of finite presentation on an
analytic space \(S\); the functor \(F_{\mathscr{A}}\) is representable
in the category \(\mathsf{An}/S\) by an analytic space separated over
\(S\).

\end{itenv}

\begin{proof}
We note first of all that the functor \(F_{\mathscr{A}}\) is of a local
nature (cf.~Exposé~11, Definition~5.4), since if \(T\) is an analytic
space over \(S\), then the presheaf \(U\mapsto F_{\mathscr{A}}(U)\)
(where \(U\) runs over the opens of \(T\)) is precisely the sheaf of
germs of algebra homomorphisms
\(\underline{\operatorname{Hom}}_{{\mathscr{O}}_T}({\mathscr{A}}_T,{\mathscr{O}}_T)\).

Thus, by Corollary~5.7 of Exposé~11, it suffices, to show that
\(F_{\mathscr{A}}\) is representable, to find an open cover \((S_i)\) of
\(S\) such that, for each index \(i\), the functor
\(F_{\mathscr{A}}/S_i\) is representable in \(\mathsf{An}/S_i\).

By \hyperref[definition-1]{Definition~1}, we thus see that it suffices
to consider the case where \({\mathscr{A}}\) is of the form
\({\mathscr{O}}_S[t_1,\ldots,t_n]/(f_1,\ldots,f_m)\). Then the exact
sequence \oldpage{19-02} \[
  0 \to {\mathscr{I}} \xrightarrow{i} {\mathscr{O}}_S[t_1,\ldots,t_n] \to {\mathscr{A}}\to 0
\] (where \({\mathscr{I}}=(f_1,\ldots,f_m)\)) gives, by base extension,
an exact sequence \[
  {\mathscr{I}}_T \xrightarrow{i_T} {\mathscr{O}}_T[t_1,\ldots,t_n] \to {\mathscr{A}}_T \to 0
\] for every analytic space \(T\) over \(S\). This allows us to identify
\(F_{\mathscr{A}}(T)=\operatorname{Hom}_{{\mathscr{O}}_T}({\mathscr{A}}_T,{\mathscr{O}}_T)\)
with the set of homomorphisms from \({\mathscr{O}}_T[t_1,\ldots,t_n]\)
to \({\mathscr{O}}_T\) that vanish on \(i_T({\mathscr{I}}_T)\).

But the functor
\(F_{{\mathscr{O}}_S[t_1,\ldots,t_n]}\colon T\mapsto\operatorname{Hom}_{{\mathscr{O}}_T}({\mathscr{O}}_T[t_1,\ldots,t_n],{\mathscr{O}}_T)\)
is isomorphic to \(T\mapsto(\Gamma(T,{\mathscr{O}}_T))^n\); by
Theorem~1.1 of Exposé~10 and Proposition~3.1 of Exposé~11, we see that
\(F_{{\mathscr{O}}_S[t_1,\ldots,t_n]}\) is represented in
\(\mathsf{An}/S\) by the pair \((S\times{\mathscr{E}}^n,\eta)\), where
\(\eta\) is the homomorphism of algebras
\(\eta\colon{\mathscr{O}}_{S\times{\mathscr{E}}^n}[t_1,\ldots,t_n]\to{\mathscr{O}}_{S\times{\mathscr{E}}^n}\)
defined by \(\eta(t_i)=q^*(z_i)\) (for \(i=1,\ldots,n\)) and denoting by
\(q\) the projection from \(S\times{\mathscr{E}}^n\) onto
\({\mathscr{E}}^n\). It immediately follows that the sub-functor
\(F_{\mathscr{A}}\) is represented by the closed subspace
\(X\subset S\times{\mathscr{E}}^n\) defined by the ideal
\({\mathscr{I}}=\eta\circ i_{S\times{\mathscr{E}}^n}({\mathscr{I}}_{S\times{\mathscr{E}}^n})\subset{\mathscr{O}}_{S\times{\mathscr{E}}^n}\)
(cf.~Exposé~11, Lemma~3.6).

In the case in question, \(X\) is evidently separated over \(S\). Since
the property of being separated over \(S\), for an object of
\(\mathsf{An}/S\), is local on \(S\), the latter claim of the statement
is also proven (cf.~Expose~11, Proposition~5.6: we obtain the space that
represents \(F_{\mathscr{A}}\) by gluing together those that represent
the \(F_{{\mathscr{A}}/S_i}\)).
\end{proof}

\protect\phantomsection\label{definition-2}
\begin{rmenv}{Definition 2}
For every \({\mathscr{O}}_S\)-algebra \({\mathscr{A}}\) of finite
presentation, we define the \emph{analytic spectrum} of
\({\mathscr{A}}\), denoted
\(\operatorname{Spec}_\mathrm{an}({\mathscr{A}})\), to be the object of
\(\mathsf{An}/S\) (defined up to \(S\)-isomorphism) that represents the
functor \[
  F_{\mathscr{A}}\colon T\longmapsto \operatorname{Hom}_{{\mathscr{O}}_T}({\mathscr{A}}_T,{\mathscr{O}}_T).
\] \oldpage{19-3} \(\operatorname{Spec}_\mathrm{an}({\mathscr{A}})\) is
thus an analytic space separated over \(S\).

\end{rmenv}

\begin{rmenv}{Example}
Consider an \({\mathscr{O}}_S\)-module \({\mathscr{E}}\) of finite
presentation; its symmetric algebra \(S({\mathscr{E}})\) is an
\({\mathscr{O}}_S\)-algebra of finite presentation, and
\(\operatorname{Spec}_\mathrm{an}(S({\mathscr{E}}))\) represents the
functor \[
  T\longmapsto
  \operatorname{Hom}_{{{\mathscr{O}}_T}\text{-}\mathsf{Alg}}(S({\mathscr{E}}_T),{\mathscr{O}}_T)
  \simeq \operatorname{Hom}_{{{\mathscr{O}}_T}\text{-}\mathsf{Mod}}({\mathscr{E}}_T,{\mathscr{O}}_T).
\] This is the \emph{vector bundle} defined by \({\mathscr{E}}\)
(cf.~Exposé~12, §1).

\end{rmenv}

\protect\phantomsection\label{proposition-2}
\begin{itenv}{Proposition 2}

\begin{enumerate}
\def\labelenumi{\roman{enumi}.}
\tightlist
\item
  \({\mathscr{A}}\mapsto\operatorname{Spec}_\mathrm{an}({\mathscr{A}})\)
  defines a contravariant functor from the category of
  \({\mathscr{O}}_S\)-algebras of finite presentation to
  \(\mathsf{An}/S\).
\item
  For any \({\mathscr{O}}_S\)-algebras \({\mathscr{A}}\) and
  \({\mathscr{B}}\) of finite presentation,
  \({\mathscr{A}}\otimes_{{\mathscr{O}}_S}{\mathscr{B}}\) is of finite
  presentation, and
  \(\operatorname{Spec}_\mathrm{an}({\mathscr{A}}\otimes_{{\mathscr{O}}_S}{\mathscr{B}})\simeq\operatorname{Spec}_\mathrm{an}({\mathscr{A}})\times_S\operatorname{Spec}_\mathrm{an}({\mathscr{B}})\).
\item
  Let \(h\colon{\mathscr{A}}\to{\mathscr{B}}\) be a surjective
  homomorphism of algebras of finite presentation. Then the
  corresponding morphism
  \(\operatorname{Spec}_\mathrm{an}(h)\colon\operatorname{Spec}_\mathrm{an}({\mathscr{B}})\to\operatorname{Spec}_\mathrm{an}({\mathscr{A}})\)
  is a closed immersion.
\item
  For every \({\mathscr{O}}_S\)-algebra \({\mathscr{A}}\) of finite
  presentation, and for every analytic space \(T\) over \(S\),
  \(\operatorname{Spec}_\mathrm{an}({\mathscr{A}})\times_S T\simeq\operatorname{Spec}_\mathrm{an}({\mathscr{A}}_T)\)
  (compatibility with base change).
\end{enumerate}

\end{itenv}

\begin{proof}
\leavevmode

\begin{enumerate}
\def\labelenumi{(\roman{enumi})}
\tightlist
\item
  is immediate, since \(F_{\mathscr{A}}(T)\) is a contravariant functor
  in \({\mathscr{A}}\);
\item
  follows from the formula \[
    \operatorname{Hom}_{{\mathscr{O}}_T}({\mathscr{A}}_T\otimes_{{\mathscr{O}}_T}{\mathscr{B}}_T,{\mathscr{O}}_T)
    \simeq \operatorname{Hom}_{{\mathscr{O}}_T}({\mathscr{A}}_T,{\mathscr{O}}_T)\times\operatorname{Hom}_{{\mathscr{O}}_T}({\mathscr{B}}_T,{\mathscr{O}}_T)
  \]
\item
  follows easily from the proof of
  \hyperref[proposition-1]{Proposition~1}, and (iv) can be proven by
  Corollary~3.2 of Exposé~11.
\end{enumerate}

\end{proof}

Consider the pair \((X,\xi)\) that represents the functor
\(F_{\mathscr{A}}\) defined by an \({\mathscr{O}}_S\)-algebra
\({\mathscr{A}}\) of finite presentation. Then
\(X=\operatorname{Spec}_\mathrm{an}({\mathscr{A}})\) is an analytic
space over \(S\), with structure morphism \(f\colon X\to S\), and
\(\xi\in F_{\mathscr{A}}(X)=\operatorname{Hom}_{{\mathscr{O}}_X}({\mathscr{A}}_X,{\mathscr{O}}_X)\)
is a homomorphism from \({\mathscr{A}}_X=f^*({\mathscr{A}})\) to
\({\mathscr{O}}_X\). In other words, \(\xi\) defines a factorisation \[
  X \xrightarrow{\hat{f}} (|S|,{\mathscr{A}}) \to S
\] of the structure morphism \(f\) of \(X\) through the ringed space
\((|S|,{\mathscr{A}})\) (where \(|S|\) denotes the underlying
topological space of \(S\)) endowed with the canonical morphism
\((|S|,{\mathscr{A}})\to S\) (consisting of the identity map on \(|S|\)
and the homomorphism \({\mathscr{O}}_S\to{\mathscr{A}}\)).

\oldpage{19-04}

For any \({\mathscr{A}}\)-module \({\mathscr{M}}\), set \[
  \widetilde{{\mathscr{M}}}
  = \hat{f}^*({\mathscr{M}})
  = f^*({\mathscr{M}})\otimes_{f^*({\mathscr{A}})}{\mathscr{O}}_X.
\] We thus define a functor
\({\mathscr{M}}\mapsto\widetilde{{\mathscr{M}}}\) from the category of
\({\mathscr{A}}\)-modules to the category of \({\mathscr{O}}_X\)-modules
(where \(X=\operatorname{Spec}_\mathrm{an}({\mathscr{A}})\)). If
\({\mathscr{M}}\) and \({\mathscr{N}}\) are \({\mathscr{A}}\)-modules
then \[
  ({\mathscr{M}}\otimes_{\mathscr{A}}{\mathscr{N}})^\sim
  = \widetilde{{\mathscr{M}}}\otimes_{{\mathscr{O}}_X}\widetilde{{\mathscr{N}}}.
\]

Consider an analytic space \(T\) over \(S\), and the space \[
  X_T
  = \operatorname{Spec}_\mathrm{an}({\mathscr{A}}_T)
  = X\times_S T
\] over \(T\) ((iv) of \hyperref[proposition-2]{Proposition~2}); the
structure morphism from \(X_T\) to \(T\) again factors as \[
  X_T
  \xrightarrow{\hat{f}_T}
  (|T|,{\mathscr{A}}_T)
  \to T
\] and we obtain a commutative diagram \[
  \begin{CD}
    X @<<< X_T
  \\@V{\hat{f}}VV @VV{\hat{f}_T}V
  \\(|S|,{\mathscr{A}}) @<<< (|T|,{\mathscr{A}}_T)
  \end{CD}
\] Let \({\mathscr{M}}\) be an \({\mathscr{A}}\)-module; the inverse
image \({\mathscr{M}}_T={\mathscr{M}}\times_S T\) of \({\mathscr{M}}\)
over \(T\) is an \({\mathscr{A}}_T\)-module. We thus see that
\(\widetilde{{\mathscr{M}}}_T=\hat{f}_T^*({\mathscr{M}}_T)\) can be
identified with the inverse image of \(\widetilde{{\mathscr{M}}}\) under
the morphism \(X_T\to X\) (compatibility of the functor
\({\mathscr{M}}\mapsto\widetilde{{\mathscr{M}}}\) with base change).

\section{Fibre of the analytic spectrum over a point in the
base}\label{section-2}

Consider an \({\mathscr{O}}_S\)-algebra \({\mathscr{A}}\) of finite
presentation, and its analytic spectrum
\(X=\operatorname{Spec}_\mathrm{an}({\mathscr{A}})\). If \(s\in S\),
recall (Exposé~10, §2) that the \emph{fibre} of \(X\) at \(s\) is the
fibre product \(X_s=X\times_S\{s\}\), where \(\{s\}\) is the analytic
space consisting of the point \(s\), endowed with the field
\({\mathscr{O}}_s/{\mathfrak{m}}_s{\mathscr{O}}_s=K(s)\simeq k\) (where
\({\mathfrak{m}}_s\) denotes the maximal ideal of the local ring
\({\mathscr{O}}_s\)). \oldpage{19-05} Then (by (iv) of
\hyperref[proposition-2]{Proposition~2}) \[
  X_s
  = X\times_S\{s\}
  \simeq \operatorname{Spec}_\mathrm{an}({\mathscr{A}}_{\{s\}})
  = \operatorname{Spec}_\mathrm{an}({\mathscr{A}}(s))
\] where
\({\mathscr{A}}(s)={\mathscr{A}}_s/{\mathfrak{m}}_s{\mathscr{A}}_s={\mathscr{A}}_s\otimes_{{\mathscr{O}}_s}k={\mathscr{A}}_{\{s\}}\);
if on a neighbourhood \(U\) of \(s\) we have that \[
  {\mathscr{A}}|U \simeq \frac{{\mathscr{O}}_U[t_1,\ldots,t_n]}{(f_1,\ldots,f_m)}
\] then \[
  {\mathscr{A}}(s) \simeq \frac{k[t_1,\ldots,t_n]}{(f_1(s),\ldots,f_m(s))}.
\] \(X_s=\operatorname{Spec}_\mathrm{an}({\mathscr{A}}(s))\) is the
closed subspace of \(\{s\}\times{\mathscr{E}}^n\simeq{\mathscr{E}}^n\)
defined by the ``equations'' \[
  f_i(s)(z_1,\ldots,z_n) = 0
  \qquad\text{for }i=1,2,\ldots,m.
\]

More precisely, if the functor \(F_{\mathscr{A}}\) is represented in
\(\mathsf{An}/S\) by the pair \((X,\xi)\), then the functor
\(F_{{\mathscr{A}}(s)}=F_{{\mathscr{A}}/\{s\}}\) is represented in
\(\mathsf{An}/\{s\}\) by the pair \((X_s,\xi_s)\), where \(X_s\) is the
fibre of \(X\) at \(s\) and where
\(\xi_s\colon{\mathscr{A}}_{X_s}\to{\mathscr{O}}_{X_s}\) is obtained
from \(\xi\colon{\mathscr{A}}_X\to{\mathscr{O}}_X\) by reduction modulo
\({\mathfrak{m}}_s\). In particular, the map from
\(\operatorname{Hom}_{\{s\}}(\{s\},X_s)\) to
\(\operatorname{Hom}_k({\mathscr{A}}(s),k)\) that sends a section \(f\)
of \(X_s/\{s\}\) to the composite homomorphism \(f^*\circ\xi'_{s,f(s)}\)
(where
\(\xi'_{s,f(s)}\colon{\mathscr{A}}(s)\to{\mathscr{O}}_{X_{s,f(s)}}\) is
induced by the germ of \(\xi_s\) at \(f(s)\)) is bijective. Since
\(\operatorname{Hom}(\{s\},X_s)\) can be identified with the set of
points of \(X_s\) by sending \(f\colon\{s\}\to X_s\) to the point
\(f(s)=x\) (so that \(f^*=\epsilon_X\colon{\mathscr{O}}_{X_s,x}\to k\)
is the augmentation homomorphism), we obtain a bijection between the set
of points of \(X_s\) and the set of ideals of \({\mathscr{A}}(s)\) that
give \(k\) as their quotient, by sending \(x\in X_s\) to the inverse
image under \(\xi'_{s,x}\colon{\mathscr{A}}(s)\to{\mathscr{O}}_{X_s,x}\)
of the maximal ideal of \({\mathscr{O}}_{X_s,x}\). But
\({\mathscr{O}}_{X_s,x}={\mathscr{O}}_{X,x}/{\mathfrak{m}}_s{\mathscr{O}}_{X,x}\)
(Exposé~10, §2),
\({\mathscr{A}}(s)={\mathscr{A}}_s/{\mathfrak{m}}_s{\mathscr{A}}_s\),
and \(\xi'_{s,x}\) is obtained by reduction modulo \({\mathfrak{m}}_s\)
of \(\xi'_X\colon{\mathscr{A}}_s\to{\mathscr{O}}_{X,x}\) (induced by the
germ of \(\xi\) at \(x\)). We can thus state:

\oldpage{19-06}

\protect\phantomsection\label{proposition-3}
\begin{itenv}{Proposition 3}
We

\end{itenv}


% Bibliography
\nocite{*}
\printbibliography[heading=bibintoc,title=Bibliography]

\end{document}
